% !TeX root = ../../MainBook.tex
\section{傅里叶反演公式}\label{sec:Approach}

\subsection{\texorpdfstring{$\delta$分布的逼近}{delta分布的逼近}}\label{sub:approach}
在\refsub{shannon}中，我们对带限函数提出了公式：如果$\supp(\mathcal{F} f)\subset\sbr{-\frac{T}{2},\frac{T}{2}}$，则有
\begin{equation}
    f(t)=T\sinc(Tt)*f(t).
\end{equation}

如果希望对更一般的函数$f(x)$，找到一个函数$g(x)$，使得$f(x)=g*f(x)$，自然会想到Dirac $\delta$分布，然而它是一个奇异分布，我们只能从形式上定义一个$\delta$函数。事实上，
\begin{equation}
    \nexists g,\quad f*g=f,\quad\forall f,
\end{equation}
但很容易构造一个函数列$g_n$，使得$f*g_n\to f$\cite{folland}。回忆我们在\refsub{intuition_of_convolution}中提到的“平均化”卷积，设
\begin{equation}
    g_a(x)=\frac{1}{2a}\chi_{[-a,a]},
\end{equation}
则$g_a*f(t)$是$f$在区间$[t-a,t+a]$上的平均值，即
\begin{equation}
    f*g_a(t)=\frac{1}{2a}\int_{t-a}^{t+a}f(x)\,\dd x.
\end{equation}
若$f\in C(\mathbb{R})$，则$g_a*f\to f,\ a\to 0$。

一般地，我们考虑函数的缩放，设$g\in L^1(\mathbb{R})$，我们记
\begin{equation}
    g_\lambda(x)=\frac{1}{\lambda}g\paren{\frac{x}{\lambda}}.
\end{equation}
这类似于时域的伸缩在傅里叶变换下的表现，它是一种“保面积”的缩放，即：
\begin{equation}
    \int_{-\infty}^{\infty}g_\lambda(x)\,\dd x=\int_{-\infty}^{\infty}g\paren{\frac{y}{\lambda}}\,\dd\frac{y}{\lambda}=\int_{-\infty}^{\infty}g(y)\,\dd y.
\end{equation}
更一般地，我们给出接下来要频繁使用的引理：
\begin{lemma}
    \begin{equation}
        \int_{a}^{b}g_\lambda(x)\,\dd x=\int_{a/\lambda}^{b/\lambda}g(y)\,\dd y.
    \end{equation}
\end{lemma}

\begin{theorem}
    设$g\in L^1(\mathbb{R}),\ \int_{-\infty}^{\infty}g(x)\,\dd x=1$，并记$\alpha=\int_{-\infty}^{0}g(x)\,\dd x,\ \beta=\int_{0}^{\infty}g(x)\,\dd x$。假设$f\in PC(\mathbb{R})$，则当$g(x)$时限或$f(x)$有界时，$f*g(x)$良定义，并且
    \begin{equation}
        \lim_{\lambda\to 0}f*g_\lambda(x)=\alpha f(x+)+\beta f(x-).
    \end{equation}
    特别地，如果$f(x)$在点$x$连续，则有$\lim_{\lambda\to 0}f*g_\lambda(x)=f(x)$；如果$f(x)$在某个闭区间$[a,b]$上连续，则以上极限是一致收敛于$f(x)$的。
\end{theorem}
\begin{proof}
    以下证明中，我们需要使用$\epsilon-\delta$语言，暂时不用$\delta$表示Dirac分布。
    \begin{equation}
        f*g_\lambda(x)-\alpha f(x+)-\beta f(x-)=\int_{-\infty}^{0}[f(x-y)-f(x+)]g_\lambda(y)\,\dd y+\int_{0}^{\infty}[f(x-y)-f(x-)]g_\lambda(y)\,\dd y.
    \end{equation}
    我们希望证明在$\lambda$充分小时，以上两个积分也充分小。以$\int_{0}^{\infty}$为例，由于$f\in PC(\mathbb{R})$，存在$\delta>0$，使得当$|y|<\delta$时，$|f(x-y)-f(x-)|<\epsilon$。于是
    \begin{equation}
        \abs{\int_{0}^{\delta}[f(x-y)-f(x-)]g_\lambda(y)\,\dd y}<\epsilon\abs{\int_{0}^{\delta}g_\lambda(y)\,\dd y}=\epsilon\abs{\int_{0}^{\delta/\lambda}g(y)\,\dd y}\leq\epsilon\int_{0}^{\infty}|g(y)|\,\dd y.
    \end{equation}
    我们用剩下的条件估计$\int_{\delta}^{\infty}[f(x-y)-f(x-)]g_\lambda(y)\,\dd y$。如果$g(x)$时限，设$\supp(g)\subset[-T,T]$，则当$\lambda<\delta/T$时，
    \begin{align}
        y>\delta/\lambda>T\implies g(y)=0&\implies\int_{\delta}^{\infty}[f(x-y)-f(x-)]g_\lambda(y)\,\dd y \nonumber \\
        &=\int_{\delta/\lambda}^{\infty}[f(x-\lambda y)-f(x-)]g(y)\,\dd y=0.
    \end{align}

    如果$f(x)$有界，设$|f(x)|<M$，则
    \begin{align}
        \abs{\int_{\delta}^{\infty}[f(x-y)-f(x-)]g_\lambda(y)\,\dd y}&\leq \int_{\delta}^{\infty}|f(x-y)-f(x-)||g_\lambda(y)|\,\dd y \nonumber \\
        &\leq 2M\int_{\delta}^{\infty}|g_\lambda(y)|\,\dd y \nonumber \\
        &=2M\int_{\delta/\lambda}^{\infty}|g(y)|\,\dd y.
    \end{align}
    $\int_{-\infty}^{\infty}|g(y)|\,\dd y<\infty$的必要条件是$\int_{\delta/\lambda}^{\infty}|g(y)|\,\dd y\to 0,\ \delta/\lambda\to\infty$，因此上式随$\lambda\to 0$趋于$0$。

    最后，如果$f(x)$在闭区间$[a,b]$上连续，则它在该区间上一致连续，以上证明中的$\delta$与$x$无关，极限是一致收敛的。
\end{proof}
一种常用的函数列是高斯函数列，这里取一个积分值为$1$的特例：
\begin{equation}
    G(x)=\frac{1}{\sqrt{\pi}}\e^{-x^2},\quad G_\lambda(x)=\frac{1}{\lambda\sqrt{\pi}}\e^{-x^2/\lambda^2}.
\end{equation}
尽管用很多其他函数都可以逼近$\delta$分布，但高斯函数列有一系列非常好的性质：
\begin{itemize}
    \item 高斯函数是偶函数，$\alpha=\beta=\displaystyle\frac{1}{2}$；
    \item 高斯函数是无限阶可导的。我们知道，卷积能够继承函数的可导性：$f*g'=f'*g=(f*g)'$，使用高斯函数做卷积，所得的函数光滑
    \item 高斯函数是Schwartz函数，即$G_\lambda\in\mathcal{S}(\mathbb{R})$
    \item 高斯函数的傅里叶变换仍是高斯函数，这使得傅里叶变换更容易计算
    \item 高斯函数$G_\lambda\in L^1(\mathbb{R})$并且是有界的，与任意$f\in PC(\mathbb{R})$卷积时，卷积积分良定义且趋向于$f$
\end{itemize}

作为以上定理的推论，可以得到多项式空间在连续函数空间中的稠密性：
\begin{theorem}[*Weierstrass第一逼近定理]\label{thm:weierstrass_1}
    设$f\in C([a,b]),\ -\infty<a<b<\infty$，则$f$是某个多项式序列在$[a,b]$上的一致极限，即：
    \begin{equation}
        \forall \epsilon>0,\ \exists \text{多项式 } P,\quad \sup_{a\leq x\leq b}|f(x)-P(x)|<\epsilon.
    \end{equation}
\end{theorem}
%在处理分析问题时，类似分段光滑函数可由三角级数一致逼近，将连续函数用多项式一致逼近是一个很强大的工具。
\begin{proof}
    将$f$做连续延拓，使得
    \begin{equation}
        f\in C(\mathbb{R}),\quad f(x)=0,\quad\forall x\in (-\infty,a-1]\cup[b+1,\infty),
    \end{equation}
    则$\lim_{\lambda\to 0}f*G_\lambda \rightrightarrows f,\ \forall x\in[a,b]$，即$f*G_\lambda$在区间$[a,b]$上一致收敛到$f$。具体来说，
    \begin{equation}
        \forall \epsilon>0,\ \exists \lambda>0,\quad \sup_{a\leq x\leq b}\abs{f(x)-\frac{1}{\lambda\sqrt{\pi}}\int_{a-1}^{b+1}\e^{-(x-y)^2/\lambda^2}f(y)\,\dd y}<\frac{\epsilon}{2}.
    \end{equation}
    考虑将右边的积分用多项式近似。我们知道，$\e^{-t^2}$可以展开为Taylor级数，并且在$\mathbb{C}$上内闭一致收敛，即
    \begin{equation}
        \sum_{n=0}^{N}\frac{(-1)^n t^{2n}}{n!}\rightrightarrows \e^{-t^2},\quad N\to\infty.
    \end{equation}
    具体来说，取$|f(x)|$的上界$M$（紧支连续函数是一致连续的，必有上界），存在$N$，使得
    \begin{equation}
        \sup_{t\in[a-1,b+1]}\abs{\e^{-t^2}-\sum_{n=0}^{N}\frac{(-1)^n t^{2n}}{n!}}<\frac{\epsilon\sqrt{\pi}}{2M(b-a+2)},
    \end{equation}
    于是
    \begin{equation}
        P(x)=\frac{1}{\lambda\sqrt{\pi}}\sum_{0}^{N}\int_{a-1}^{b+1}\frac{(-1)^n (x-y)^{2n}}{\lambda^{2n}n!}f(y)\,\dd y,\quad \sup_{a\leq x\leq b}|f(x)-P(x)|<\epsilon.
    \end{equation}
    我们还可以求出$P$的显式表达式：
    \begin{equation}
        P(x)=\sum_{n=0}^{N}\sum_{k=0}^{2n}c_{kn}x^k,\quad c_{kn}=\frac{(-1)^{k-n}}{\lambda^{2n+1}n!\sqrt{\pi}}\binom{2n}{k}\int_{a-1}^{b+1}y^{2n-k}f(y)\,\dd y.
    \end{equation}
    使用其他光滑函数列逼近$\delta$分布也可以得到不同的多项式逼近，不过，绝大多数情况下我们并不需要计算出这个表达式，而是直接用多项式的一致逼近来获取函数的性质。
\end{proof}

进一步，可以得到三角多项式（\cref{def:trinomial}）空间在连续的周期函数空间中的稠密性：
\begin{corollary}[Weierstrass第二逼近定理]\label{cor:weierstrass_2}
    设$f\in C(\mathbb{R})$以$T$为周期，则$f$是某个同周期的三角多项式序列的一致极限，即：
    \begin{equation}
        \forall \epsilon>0,\ \exists \text{三角多项式 } P,\quad \sup_{x}|f(x)-P(x)|<\epsilon.
    \end{equation}
\end{corollary}
\begin{proof}
    由于可以对自变量进行伸缩变换，我们不妨设$T=2\pi$，并重点考虑区间$[-\pi,\pi]$。

    首先，三角多项式的奇分量和偶分量是显式地分离的，这启发我们考虑将$f(t)$也分解为奇分量和偶分量：
    \begin{equation}
        f_{\mathrm{e}}(t)=\frac{f(t)+f(-t)}{2},\quad f_{\mathrm{o}}(t)=\frac{f(t)-f(-t)}{2}.
    \end{equation}
    一种直接的想法是，通过变量代换$x=\cos t$和$x=\sin t$，分别将逼近$f_{\mathrm{e}}(\arccos x)$和$f_{\mathrm{o}}(\arcsin x)$的多项式转化为三角多项式。然而，$\sin(k x),\ k\in\mathbb{N}$有共同的零点和对称轴：
    \begin{equation}
        \sin(-k\pi)=\sin 0=\sin(k\pi)=0,\quad \sin(k(\pi-x))=\sin(kx)=\sin(k(-\pi-x)),
    \end{equation}
    对于奇分量来说，我们只能保证
    \begin{equation}
        f_{\mathrm{o}}(-\pi)=f_{\mathrm{o}}(0)=f_{\mathrm{o}}(\pi)=0,
    \end{equation}
    而不能保证$f_{\mathrm{o}}(t)$关于$t=\pm\displaystyle\frac{\pi}{2}$轴对称，因此不能单独使用正弦多项式逼近奇分量。

    考虑以下两个偶函数：
    \begin{equation}
        g(t)=f(t)+f(-t),\quad h(t)=\sin(t)(f(t)-f(-t)).
    \end{equation}
    做变量代换$x=\cos t$，利用\cref{thm:weierstrass_1}，有：
    \begin{equation}
        \forall \epsilon>0,\ \exists \text{多项式 } P,Q,\quad \sup_{-1\leq x\leq 1}|g(\arccos x)-P(x)|<\epsilon,\quad \sup_{-1\leq x\leq 1}|h(\arccos x)-Q(x)|<\epsilon.
    \end{equation}
    代回变量$t$，$P(\cos t)$和$Q(\cos t)$都是以$2\pi$为周期的偶函数，符合我们的要求。

    我们希望能用三角多项式逼近$f(t)-f(-t)$，也就是说$Q(\cos t)/\sin t$仍为三角多项式，这实际上要求$Q(\cos t)$在$t=0$和$\pm\pi$处，即$x=\pm 1$处有零点。注意到$x=\pm 1$时，有：
    \begin{equation}
        h(\arccos x)=0,\quad |Q(x)-h(\arccos x)|=|Q(x)|<\epsilon.
    \end{equation}
    构造修正多项式
    \begin{equation}
        R(x)=Q(x)-Q(1)\frac{1+x}{2}-Q(-1)\frac{1-x}{2}.
    \end{equation}
    这样，$R(\pm 1)=0$，且
    \begin{align}
        |R(x)-Q(x)|&=\abs{Q(1)\frac{1+x}{2}+Q(-1)\frac{1-x}{2}} \nonumber \\
        &\leq |Q(1)|\frac{|1+x|}{2}+|Q(-1)|\frac{|1-x|}{2} \nonumber \\
        &\leq 2\epsilon,\quad \forall x\in[-1,1].
    \end{align}
    因此
    \begin{equation}
        \sup_{-1\leq x\leq 1}|R(x)-h(\arccos x)|\leq\sup_{-1\leq x\leq 1}|R(x)-Q(x)|+\sup_{-1\leq x\leq 1}|Q(x)-h(\arccos x)|\leq 3\epsilon.
    \end{equation}
    由于$R(x)$有零点$\pm 1$，它一定有因式$(1-x)(1+x)=1-x^2$，记$R(x)=R_0(x)(1-x^2)$，则
    \begin{equation}
        R(\cos t)=R_0(\cos t)(1-\cos^2 t)=R_0(\cos t)\sin^2 t.
    \end{equation}

    最后，我们通过分段估计来说明三角多项式$R_0(\cos t)\sin t$能够一致逼近$F(t)=f(t)-f(-t)$。一方面，由于$F(t)$连续且$F(0)=F(\pi)=0$（根据周期性，$F(-\pi)$处的情况与$F(\pi)$处相同），有
    \begin{equation}
        \forall\eta>0,\ \exists\delta_1>0,\quad |t|<\delta_1\ \text{或}\ |t-\pi|<\delta_1\implies |F(t)|<\frac{\eta}{2}.
    \end{equation}
    取$R_0(\cos t)$的上界$M$，根据$\sin t$的连续性，
    \begin{equation}
        \exists\delta_2>0,\quad |t|<\delta_2\ \text{或}\ |t-\pi|<\delta_2\implies|\sin t|<\frac{\eta}{2M}.
    \end{equation}
    令$\delta=\min\{\delta_1,\delta_2\}$，则
    \begin{equation}
        |t|<\delta\ \text{或}\ |t-\pi|<\delta\implies |F(t)-R_0(\cos t)\sin t|\leq |F(t)|+M|\sin t|\leq\eta.
    \end{equation}
    另一方面，当$|t|>\delta$且$|t-\pi|>\delta$时，$|\sin t|>\sin\delta>0$，因此
    \begin{equation}
        |F(t)\sin t-R_0(\cos t)\sin^2 t|<\epsilon\implies|F(t)-R_0(\cos t)\sin t|<\frac{\epsilon}{\sin\delta}.
    \end{equation}
    只要选取$\epsilon<\eta\sin\delta$，就能保证$\sup_t |F(t)-R_0(\cos t)\sin t|<\eta$。
\end{proof}

\subsection{Riemann-Lebesgue引理}\label{sub:riemman_lebesgue_lemma}
本小节来补充\refsub{introduction_of_fourier}中提到的\textbf{Riemann-Lebesgue引理}的证明。
\begin{theorem}[Riemann-Lebesgue引理]\label{thm:riemann_lebesgue}
    设$f\in L^1(\mathbb{R})$，则其傅里叶变换$\mathcal{F} f(x)\to 0,\ x\to\infty$。
\end{theorem}
我们给出一个基于阶梯函数逼近的证明，尽管现在不需要掌握它，但可以从这个证明体会到分析学中逼近技术的强大威力。
\begin{proof}
    首先假设$f$是阶梯函数，即$f(t)=\sum_{k=1}^{n}c_k \chi_{[a_k,b_k]}(t)$，我们知道：
    \begin{lralign}
        \mathcal{F} \chi_{[a,b]}(\xi)=\int_{a}^{b}\e^{-2\pi\im\xi t}\,\dd t=\frac{\e^{-2\pi\im\xi a}-\e^{-2\pi\im\xi b}}{2\pi\im\xi}
    \switch
        \mathcal{F} \chi_{[a,b]}(\omega)=\int_{a}^{b}\e^{-\im\omega t}\,\dd t=\frac{\e^{-\im\omega a}-\e^{-\im\omega b}}{\im\omega}.
    \end{lralign}
    分子是两个模值为$1$的复数相减，分母趋于无穷，故
    \begin{equation}
        \mathcal{F} \chi_{[a,b]}(x)\to 0,\quad x\to\infty.
    \end{equation}
    由线性性，得到：
    \begin{equation}
        \mathcal{F} f(x)=\sum_{k=1}^{n}c_k\mathcal{F}\chi_{[a_k,b_k]}\to 0,\quad x\to\infty.
    \end{equation}

    设$f\in L^1(\mathbb{R})$，则根据实变函数的理论，存在阶梯函数列$f_n$，使得$\|f-f_n\|_1\to 0,\ n\to\infty$。我们知道，
    \begin{equation}
        |\mathcal{F} g(x)|\leq \int_{-\infty}^{\infty}|g(t)|\,\dd t=\|g\|_1<\infty,\quad g\in L^1(\mathbb{R}),
    \end{equation}
    因此，
    \begin{equation}
        \sup_{x\in\mathbb{R}}|\mathcal{F} f(x)-\mathcal{F} f_n(x)|\leq \|f-f_n\|_1\to 0,\quad n\to\infty,
    \end{equation}
    即$\mathcal{F} f_n\rightrightarrows\mathcal{F} f,\ \forall x\in\mathbb{R}$，于是$\mathcal{F} f(x)\to 0,\ x\to\infty$，得证。
\end{proof}

\subsection{傅里叶反演公式}\label{sub:fourier_inversion}
\begin{theorem}[傅里叶反演公式]
    \label{thm:fourier_inversion_1}
    设$f\in L^1(\mathbb{R})\cap PC(\mathbb{R})$，则对于任意点$x$，都有
    \begin{lralign}
        \frac{f(x+)+f(x-)}{2}&=\lim_{\lambda\to 0}\int_{-\infty}^{\infty}\e^{2\pi\im\xi x}\mathcal{F} f(\xi)\e^{-\pi^2\lambda^2\xi^2}\,\dd\xi
    \switch
        \frac{f(x+)+f(x-)}{2}&=\lim_{\lambda\to 0}\frac{1}{2\pi}\int_{-\infty}^{\infty}\e^{\im\omega x}\mathcal{F} f(\omega)\e^{-\lambda^2\omega^2/4}\,\dd\omega,
    \end{lralign}
    如果还有$\mathcal{F} f(x)\in L^1(\mathbb{R})$，则有
    \begin{lralign}
        \frac{f(x+)+f(x-)}{2}&=\int_{-\infty}^{\infty}\e^{2\pi\im\xi x}\mathcal{F} f(\xi)\,\dd\xi
    \switch
        \frac{f(x+)+f(x-)}{2}&=\frac{1}{2\pi}\int_{-\infty}^{\infty}\e^{\im\omega x}\mathcal{F} f(\omega)\,\dd\omega.
    \end{lralign}
\end{theorem}
我们指出，$\mathcal{S}(\mathbb{R})\subset L^1(\mathbb{R})\cap PC(\mathbb{R})$。

要说明傅里叶逆变换的确能够还原出时域函数，一种直接的想法是积分换序，即
\begin{lralign}
    &\int_{-\infty}^{\infty}\e^{2\pi\im\xi x}\mathcal{F} f(\xi)\,\dd\xi\\
    &=\int_{-\infty}^{\infty}\paren{\int_{-\infty}^{\infty}f(y)\e^{-2\pi\im\xi y}\,\dd y}\e^{2\pi\im\xi x}\,\dd\xi
\switch
    &\frac{1}{2\pi}\int_{-\infty}^{\infty}\e^{\im\omega x}\mathcal{F} f(\omega)\,\dd\omega\\
    &=\frac{1}{2\pi}\int_{-\infty}^{\infty}\paren{\int_{-\infty}^{\infty}f(y)\e^{-\im\omega y}\,\dd y}\e^{\im\omega x}\,\dd\omega.
\end{lralign}
然而，
\begin{lralign}
    \int_{-\infty}^{\infty}\e^{2\pi\im\xi (x-y)}\,\dd\xi
\switch
    \int_{-\infty}^{\infty}\e^{\im\omega (x-y)}\,\dd\omega
\end{lralign}
两个积分是发散的，无法进行积分换序。我们可以使用“收敛因子”来处理这个问题，从信号处理的角度，这就是\textbf{加窗}(windowing)的手法，见\refsub{window}。这里我们用高斯窗，仍记$G_\lambda(\xi)=\displaystyle\frac{1}{\lambda\sqrt{\pi}}\e^{-\xi^2/\lambda^2}$，有
\begin{lralign}
    \mathcal{F} G_\lambda(\xi)=\e^{-\pi^2\lambda^2\xi^2}
\switch
    \mathcal{F} G_\lambda(\omega)=\e^{-\lambda^2\omega^2/4}.
\end{lralign}
下面的证明中将用到高斯函数的傅里叶逆变换，但我们可以不依赖于傅里叶反演公式而预先求出它，具体的做法可以参考\cref{exmp:gauss_fourier}中求傅里叶正变换的过程。
\begin{proof}
    考虑
    \begin{lralign}
        &\int_{-\infty}^{\infty}\e^{2\pi\im\xi x}\mathcal{F} f(\xi)\e^{-\pi^2\lambda^2\xi^2}\,\dd\xi\\
        &=\int_{-\infty}^{\infty}\paren{\int_{-\infty}^{\infty}f(y)\e^{-2\pi\im\xi y}\,\dd y}\e^{2\pi\im\xi x}\e^{-\pi^2\lambda^2\xi^2}\,\dd\xi\\
        &=\int_{-\infty}^{\infty}f(y)\,\dd y\int_{-\infty}^{\infty}\e^{2\pi\im\xi (x-y)}\e^{-\pi^2\lambda^2\xi^2}\,\dd\xi
    \switch
        &\frac{1}{2\pi}\int_{-\infty}^{\infty}\e^{\im\omega x}\mathcal{F} f(\omega)\e^{-\lambda^2\omega^2/4}\,\dd\omega\\
        &=\frac{1}{2\pi}\int_{-\infty}^{\infty}\paren{\int_{-\infty}^{\infty}f(y)\e^{-\im\omega y}\,\dd y}\e^{\im\omega x}\e^{-\lambda^2\omega^2/4}\,\dd\omega\\
        &=\frac{1}{2\pi}\int_{-\infty}^{\infty}f(y)\,\dd y\int_{-\infty}^{\infty}\e^{\im\omega (x-y)}\e^{-\lambda^2\omega^2/4}\,\dd\omega.
    \end{lralign}

    以上过程利用了高斯函数的速降性质，保证了积分的绝对收敛性，从而可以进行积分换序。内层的积分恰为傅里叶变换：
    \begin{lralign}
        &\int_{-\infty}^{\infty}\e^{2\pi\im\xi (x-y)}\e^{-\pi^2\lambda^2\xi^2}\,\dd\xi\\
        &=\mathcal{F}^{-1}\sbr{\e^{-\pi^2\lambda^2\xi^2}}(x-y)\\
        &=\frac{1}{\sqrt{\pi}\lambda}\e^{-(x-y)^2/\lambda^2}\\
        &=G_{\lambda}(x-y)
    \switch
        &\frac{1}{2\pi}\int_{-\infty}^{\infty}\e^{\im\omega (x-y)}\e^{-\lambda^2\omega^2/4}\,\dd\omega\\
        &=\mathcal{F}^{-1}\sbr{\e^{-\lambda^2\omega^2/4}}(x-y)\\
        &=\frac{1}{\sqrt{\pi}\lambda}\e^{-(x-y)^2/\lambda^2}\\
        &=G_{\lambda}(x-y).
    \end{lralign}
    于是我们得到
    \begin{lralign}
        &\int_{-\infty}^{\infty}\e^{2\pi\im\xi x}\mathcal{F} f(\xi)\e^{-\pi^2\lambda^2\xi^2}\,\dd\xi\\
        &=\int_{-\infty}^{\infty}f(y)G_{\lambda}(x-y)\,\dd y\\
        &=f*G_{\lambda}(x)
    \switch
        &\frac{1}{2\pi}\int_{-\infty}^{\infty}\e^{\im\omega x}\mathcal{F} f(\omega)\e^{-\lambda^2\omega^2/4}\,\dd\omega\\
        &=\int_{-\infty}^{\infty}f(y)G_{\lambda}(x-y)\,\dd y\\
        &=f*G_{\lambda}(x).
    \end{lralign}
    我们知道，$\lim_{\lambda\to 0}f*G_{\lambda}(x)=\displaystyle\frac{f(x+)+f(x-)}{2}$，因此
    \begin{lralign}
        \lim_{\lambda\to 0}\int_{-\infty}^{\infty}\e^{2\pi\im\xi x}\mathcal{F} f(\xi)\e^{-\pi^2\lambda^2\xi^2}\,\dd\xi&=\frac{f(x+)+f(x-)}{2}
    \switch
        \lim_{\lambda\to 0}\frac{1}{2\pi}\int_{-\infty}^{\infty}\e^{\im\omega x}\mathcal{F} f(\omega)\e^{-\lambda^2\omega^2/4}\,\dd\omega&=\frac{f(x+)+f(x-)}{2}.
    \end{lralign}
    现在，只需要证明$\lim_{\lambda\to 0}$与积分号可以交换顺序。我们有：
    \begin{lralign}
        \abs{\e^{2\pi\im\xi x}\mathcal{F} f(\xi)\e^{-\pi^2\lambda^2\xi^2}}&\leq |\mathcal{F} f(\xi)|\in L^1(\mathbb{R})
    \switch
        \abs{\frac{1}{2\pi}\e^{\im\omega x}\mathcal{F} f(\omega)\e^{-\lambda^2\omega^2/4}}&\leq \frac{1}{2\pi}|\mathcal{F} f(\omega)|\in L^1(\mathbb{R}),
    \end{lralign}
    因此可以使用控制收敛定理，得到
    \begin{lralign}
        \frac{f(x+)+f(x-)}{2}&=\lim_{\lambda\to 0}\int_{-\infty}^{\infty}\e^{2\pi\im\xi x}\mathcal{F} f(\xi)\e^{-\pi^2\lambda^2\xi^2}\,\dd\xi\\
        &=\int_{-\infty}^{\infty}\e^{2\pi\im\xi x}\mathcal{F} f(\xi)\,\dd\xi
    \switch
        \frac{f(x+)+f(x-)}{2}&=\lim_{\lambda\to 0}\frac{1}{2\pi}\int_{-\infty}^{\infty}\e^{\im\omega x}\mathcal{F} f(\omega)\e^{-\lambda^2\omega^2/4}\,\dd\omega\\
        &=\frac{1}{2\pi}\int_{-\infty}^{\infty}\e^{\im\omega x}\mathcal{F} f(\omega)\,\dd\omega.
    \end{lralign}
\end{proof}
傅里叶反演公式还有一些其他的表述。例如，如果我们将窗函数换为矩形窗，则可得到：
\begin{theorem}[傅里叶反演公式]
    \label{thm:fourier_inversion}
    设$f\in L^1(\mathbb{R})\cap PS(\mathbb{R})$，则对于任意点$x$，都有
    \begin{lralign}
        \frac{f(x+)+f(x-)}{2}&=\lim_{R\to\infty}\int_{-R}^{R}\e^{2\pi\im\xi x}\mathcal{F} f(\xi)\,\dd\xi
    \switch
        \frac{f(x+)+f(x-)}{2}&=\lim_{R\to\infty}\frac{1}{2\pi}\int_{-R}^{R}\e^{\im\omega x}\mathcal{F} f(\omega)\,\dd\omega.
    \end{lralign}
\end{theorem}
等式右侧的积分与反常积分是有区别的，无穷限的反常积分要求积分上下限以二重极限的方式趋于无穷，而这里给出的是\textbf{主值意义下的积分}，即
\begin{equation}
    \int_{-\infty}^{\infty}f(x)\,\dd x=\lim_{a\to-\infty,\ b\to\infty}\int_{a}^{b}f(x)\,\dd x\neq \lim_{R\to\infty}\int_{-R}^{R}f(x)\,\dd x=\pv\int_{-\infty}^{\infty}f(x)\,\dd x.
\end{equation}
\begin{example}
    设$f(x)=x$，则$\int_{-\infty}^{\infty}f(x)\,\dd x$发散，但$\pv\int_{-\infty}^{\infty}f(x)\,\dd x=0$。
\end{example}
\begin{proof}
    我们有
    \begin{lralign}
        \int_{-R}^{R}\e^{2\pi\im\xi x}\,\dd\xi&=\frac{\e^{2\pi\im R x}-\e^{-2\pi\im R x}}{2\pi\im x}\\
        &=\frac{\sin(2\pi R x)}{\pi x}
    \switch
        \frac{1}{2\pi}\int_{-R}^{R}\e^{\im\omega x}\,\dd\omega&=\frac{\e^{\im R x}-\e^{-\im R x}}{2\im x}\\
        &=\frac{\sin(R x)}{\pi x},
    \end{lralign}
    以及\cref{thm:dirichlet_int}：
    \begin{equation}
        \int_{-\infty}^{0}\frac{\sin(ax)}{\pi x}\,\dd x=\int_{0}^{\infty}\frac{\sin(ax)}{\pi x}\,\dd x=\frac{1}{2}.
    \end{equation}
    因此
    \begin{lralign}
        \int_{-R}^{R}\e^{2\pi\im\xi x}\mathcal{F} f(\xi)\,\dd\xi&=\int_{-\infty}^{\infty}f(y)\,\dd y\int_{-R}^{R}\e^{2\pi\im\xi (x-y)}\,\dd\xi\\
        &=\int_{-\infty}^{\infty}f(y)\frac{\sin(2\pi R (x-y))}{\pi (x-y)}\,\dd y\\
        &=f*\paren{\frac{\sin(2\pi R x)}{\pi x}}
    \switch
        \frac{1}{2\pi}\int_{-R}^{R}\e^{\im\omega x}\mathcal{F} f(\omega)\,\dd\omega&=\int_{-\infty}^{\infty}f(y)\,\dd y\frac{1}{2\pi}\int_{-R}^{R}\e^{\im\omega (x-y)}\,\dd\omega\\
        &=\int_{-\infty}^{\infty}f(y)\frac{\sin(R (x-y))}{\pi (x-y)}\,\dd y\\
        &=f*\paren{\frac{\sin(R x)}{\pi x}}.
    \end{lralign}
    但是，由于$\sinc\notin L^1(\mathbb{R})$，无法保证$\lim_{R\to\infty}f(x)*\sinc(R x)=f(x)$，因此需要条件$f\in L^1(\mathbb{R})\cap PS(\mathbb{R})$，做更细致的估计：
    \begin{lralign}
        &\int_{-R}^{R}\e^{2\pi\im\xi x}\mathcal{F} f(\xi)\,\dd\xi-\frac{f(x+)+f(x-)}{2}\\
        &=\int_{-\infty}^{\infty}f(y)\frac{\sin(2\pi R (x-y))}{\pi (x-y)}\,\dd y-\frac{f(x+)+f(x-)}{2}\\
        &=\int_{-\infty}^{\infty}f(x-y)\frac{\sin(2\pi R y)}{\pi y}\,\dd y-\frac{f(x+)+f(x-)}{2}\\
        &=\int_{-\infty}^{0}[f(x-y)-f(x+)]\frac{\sin(2\pi R y)}{\pi y}\,\dd y\\
        &\quad+\int_{0}^{\infty}[f(x-y)-f(x-)]\frac{\sin(2\pi R y)}{\pi y}\,\dd y
    \switch
        &\frac{1}{2\pi}\int_{-R}^{R}\e^{\im\omega x}\mathcal{F} f(\omega)\,\dd\omega-\frac{f(x+)+f(x-)}{2}\\
        &=\int_{-\infty}^{\infty}f(y)\frac{\sin(R (x-y))}{\pi (x-y)}\,\dd y-\frac{f(x+)+f(x-)}{2}\\
        &=\int_{-\infty}^{\infty}f(x-y)\frac{\sin(R y)}{\pi y}\,\dd y-\frac{f(x+)+f(x-)}{2}\\
        &=\int_{-\infty}^{0}[f(x-y)-f(x+)]\frac{\sin(R y)}{\pi y}\,\dd y\\
        &\quad+\int_{0}^{\infty}[f(x-y)-f(x-)]\frac{\sin(R y)}{\pi y}\,\dd y.
    \end{lralign}
    我们需要证明当$R\to\infty$时，以上两个积分都趋于$0$。以$\int_{0}^{\infty}$为例，仍然考虑分段估计：$\int_{0}^{\infty}=\int_{0}^{K}+\int_{K}^{\infty}$，取$K\geq 1$，有
    \begin{lralign}
        \abs{\int_{K}^{\infty}\frac{\sin(2\pi Ry)}{\pi y}f(x-y)\,\dd y}&\leq \int_{K}^{\infty}\abs{\frac{\sin(2\pi Ry)}{\pi y}}|f(x-y)|\,\dd y\\
        &\leq \int_{K}^{\infty}|f(x-y)|\,\dd y
    \switch
        \abs{\int_{K}^{\infty}\frac{\sin(R y)}{\pi y}f(x-y)\,\dd y}&\leq \int_{K}^{\infty}\abs{\frac{\sin(R y)}{\pi y}}|f(x-y)|\,\dd y\\
        &\leq \int_{K}^{\infty}|f(x-y)|\,\dd y,
    \end{lralign}
    与
    \begin{lralign}
        \int_{K}^{\infty}\frac{\sin(2\pi Ry)}{\pi y}f(x-)\,\dd y&=f(x-)\int_{RK}^{\infty}\frac{\sin(2\pi Ry)}{\pi y}\,\dd y
    \switch
        \int_{K}^{\infty}\frac{\sin(R y)}{\pi y}f(x-)\,\dd y&=f(x-)\int_{RK}^{\infty}\frac{\sin(R y)}{\pi y}\,\dd y.
    \end{lralign}
    取$K$充分大，则以上两积分都是收敛的无穷限积分的尾部，从而是趋于$0$的。

    对于$\int_{0}^{K}$部分，定义
    \begin{equation}
        g(y)=\begin{dcases}
            \frac{f(x-y)-f(x-)}{\pi y}, & y\in[0,K], \\
            0,                          & \otherwise,
        \end{dcases}
    \end{equation}
    则
    \begin{lralign}
        &\int_{0}^{K}[f(x-y)-f(x-)]\frac{\sin(2\pi R y)}{\pi y}\,\dd y\\
        &=\int_{0}^{K}g(y)\frac{\e^{2\pi\im Ry}-\e^{-2\pi\im R y}}{2\im}\,\dd y\\
        &=\frac{1}{2\im}\sbr{\mathcal{F} g(-R)-\mathcal{F} g(R)}
    \switch
        &\int_{0}^{K}[f(x-y)-f(x-)]\frac{\sin(R y)}{\pi y}\,\dd y\\
        &=\int_{0}^{K}g(y)\frac{\e^{\im R y}-\e^{-\im R y}}{2\im}\,\dd y\\
        &=\frac{1}{2\im}\sbr{\mathcal{F} g(-R)-\mathcal{F} g(R)}.
    \end{lralign}
    由于$f\in PS(\mathbb{R}),\ g\in PC(\mathbb{R}^*),\ g(0+)=f'(x-)$，有$g\in L^1(\mathbb{R})$，根据Riemann-Lebesgue引理（\cref{thm:riemann_lebesgue}），$\mathcal{F} g(\pm R)\to 0,\ R\to\infty$，故
    \begin{equation}
        \int_{0}^{K}[f(x-y)-f(x-)]\frac{\sin(R y)}{\pi y}\,\dd y\to 0,\quad R\to\infty,
    \end{equation}
    得证。
\end{proof}

此外，还有一种类似的的傅里叶反演公式，我们仅仅给出其表述\cite{folland}：
\begin{claim}[傅里叶反演公式]
    设$f\in L^1(\mathbb{R})$，则
    \begin{lralign}
        \frac{f(x+)+f(x-)}{2}&=\lim_{\lambda\to 0}\int_{-\infty}^{\infty}\e^{2\pi\im\xi x}\mathcal{F} f(\xi)\e^{-\pi^2\lambda^2\xi^2}\,\dd\xi,\quad\aee
    \switch
        \frac{f(x+)+f(x-)}{2}&=\lim_{\lambda\to 0}\frac{1}{2\pi}\int_{-\infty}^{\infty}\e^{\im\omega x}\mathcal{F} f(\omega)\e^{-\lambda^2\omega^2/4}\,\dd\omega,\quad\aee.
    \end{lralign}
\end{claim}
